\clearpage
\phantomsection
\chapter*{Висновок}
\addcontentsline{toc}{chapter}{Висновок}
\label{chapter:conclusion}

Побудований векторний індекс у Qdrant забезпечив високу повноту пошуку за корпусом знань GTNH: показники склали \textbf{Recall@5 = 93,33\%} та \textbf{Recall@10 = 100,00\%}, що підтверджує придатність підготовлених чанків і метаданих для подальшої генерації відповідей у RAG-системі.

Порівняння ембединг-моделей показало, що для однорідного англомовного корпусу спеціалізована компактна модель \texttt{BAAI/bge-base-en-v1.5} є раціональнішим вибором, ніж важка багатомовна \texttt{BGE-m3}, оскільки дає таку саму повноту пошуку при меншому споживанні ресурсів. На корпусі обсягом близько 10 тисяч векторів параметри HNSW-індексу вже не обмежують якість видачі, а поодинокі промахи в топ-5 зумовлені непрямими згадками правил у квестах або розташуванням ключової деталі наприкінці довгого чанку -- їх доцільно усувати на наступних етапах через переранжування (reranker) або гібридний пошук.